\documentclass[conference]{IEEEtran}
\IEEEoverridecommandlockouts

\usepackage{cite}
\usepackage{amsmath,amssymb,amsfonts}
\usepackage{algorithmic}
\usepackage{graphicx}
\usepackage{textcomp}
\usepackage{xcolor}
\usepackage{booktabs}
\usepackage{url}
\usepackage{hyperref}
\hypersetup{colorlinks=true, linkcolor=black, citecolor=black, urlcolor=blue}

\def\BibTeX{{\rm B\kern-.05em{\sc i\kern-.025em b}\kern-.08em
    T\kern-.1667em\lower.7ex\hbox{E}\kern-.125emX}}

\begin{document}

\title{Distributed Fare Prediction for NYC Yellow Taxi Trips:\\
A Full-Year 2024 Analysis Using Apache Spark and HDFS}

\author{
  \IEEEauthorblockN{Kamel Al-Saweer\IEEEauthorrefmark{1},
                    Member 2\IEEEauthorrefmark{1},
                    Member 3\IEEEauthorrefmark{1},
                    Member 4\IEEEauthorrefmark{1}}
  \IEEEauthorblockA{\IEEEauthorrefmark{1}Department of Computer Science,
  Nile University, Cairo, Egypt\\
  Course: CSCI461 -- Introduction to Big Data, Spring 2026}
}

\maketitle

% ──────────────────────────────────────────────────────────────
\begin{abstract}
Predicting taxi fares accurately has significant practical value
for passengers, drivers, and urban planners. Existing studies
typically operate on sampled or single-month subsets due to the
constraints of in-memory processing frameworks. This paper presents
a fully distributed pipeline for NYC Yellow Taxi fare prediction
using the complete 12-month 2024 dataset comprising
\textbf{41,169,720 raw trip records} (660.90\,MB) published by the
NYC Taxi and Limousine Commission (TLC). The pipeline is built on a
nine-container Docker cluster running Apache Hadoop HDFS\,3.2.1 for
distributed storage and Apache Spark\,3.0 for distributed
preprocessing, exploratory data analysis (EDA), and machine
learning. After five stages of distributed cleaning, the dataset is
reduced to \textbf{35,602,215 records} with 25 engineered features.
A Spark MLlib Linear Regression model trained on 80\% of the
cleaned data achieves a Mean Absolute Error (MAE) of
\textbf{2.1607}, Root Mean Square Error (RMSE) of \textbf{5.9255},
and R$^{2}$ of \textbf{0.8934} on the held-out test set, confirming
that trip distance and duration are strong predictors of fare even
at full dataset scale. All source code and notebooks are publicly
available at \url{https://github.com/kamel1314/NYC-Taxi-BigData}.
\end{abstract}

\begin{IEEEkeywords}
Big Data, Apache Spark, HDFS, PySpark, MLlib, Linear Regression,
NYC Taxi, Fare Prediction, Distributed Computing, Docker
\end{IEEEkeywords}

% ──────────────────────────────────────────────────────────────
\section{Introduction}

Urban mobility generates massive volumes of trip data every day.
New York City's yellow taxi fleet alone recorded more than 41
million trips in 2024, each tagged with pickup and drop-off
coordinates, timestamps, passenger counts, and metered fares. This
wealth of data offers an opportunity to build accurate, data-driven
fare prediction models that benefit multiple stakeholders: passengers
can estimate costs before boarding, dispatchers can detect metering
anomalies, and city planners can identify demand patterns across
boroughs and time periods.

However, processing a full year of NYC taxi data at once poses a
genuine Big Data challenge. A single monthly file exceeds 50\,MB in
Parquet format; all twelve combined reach 660.90\,MB and contain
over 41 million rows. Conventional single-node tools such as
\textit{pandas} and \textit{scikit-learn} struggle with this volume,
either running out of memory during a join or taking hours to
complete a simple aggregation.

This paper addresses the scalability problem by building a fully
distributed pipeline on Apache Spark~\cite{zaharia2016apache} and
Hadoop HDFS~\cite{white2015hadoop}, containerised with Docker
Compose for reproducibility. The specific contributions of this work
are:

\begin{itemize}
  \item A reproducible nine-container Docker cluster integrating
        HDFS\,3.2.1, Spark\,3.0, and Apache Zeppelin\,0.10.1.
  \item A five-stage distributed preprocessing pipeline that
        reduces 41.2M raw records to 35.6M clean records with 25
        features, all executed in parallel across the cluster.
  \item A SparkSQL and Window-Function EDA revealing hour-of-day,
        day-of-week, and borough-level demand patterns.
  \item A Spark MLlib Linear Regression model achieving
        R$^{2}$\,=\,0.8934 on the full cleaned dataset, surpassing
        results reported on sampled subsets in the literature.
  \item Full public release of all code, notebooks, and
        documentation on GitHub.
\end{itemize}

The remainder of the paper is structured as follows.
Section~\ref{sec:related} reviews related work.
Section~\ref{sec:dataset} describes the dataset.
Section~\ref{sec:arch} presents the system architecture.
Section~\ref{sec:method} details the methodology.
Section~\ref{sec:results} reports experimental results.
Section~\ref{sec:discussion} discusses findings and limitations.
Section~\ref{sec:conclusion} concludes the paper.

% ──────────────────────────────────────────────────────────────
\section{Related Work}
\label{sec:related}

Fare and trip-duration prediction for taxi and ride-hailing services
has attracted considerable research attention. Huang~\cite{huang2023taxi}
compared Linear Regression, Random Forest, and XGBoost for taxi fare
prediction and found that ensemble methods consistently outperform
simple regression; however, the study used a small, single-city
dataset processed entirely in memory, limiting scalability.
Poongodi et al.~\cite{poongodi2022nyc} applied Multilayer Perceptron
(MLP) and XGBoost to the NYC TLC trip-duration task on a sampled
subset of TLC records, reporting strong accuracy (R$^{2}$ up to
0.91) while noting that single-node processing becomes a bottleneck
beyond a few million rows.

The distributed computing foundations underpinning our work trace
back to the MapReduce paradigm introduced by Dean and
Ghemawat~\cite{dean2008mapreduce} and the Google File System
architecture~\cite{ghemawat2003gfs} that inspired HDFS. Zaharia et
al.~\cite{zaharia2016apache} later showed that Spark's in-memory
directed acyclic graph (DAG) execution model reduces latency by
10--100$\times$ relative to disk-based MapReduce for iterative
workloads such as machine learning.

Our work differs from prior taxi-prediction studies in three
important ways. First, we process the \emph{complete} 12-month 2024
TLC dataset rather than a sample, operating on 41.2M rows
end-to-end. Second, model training itself is distributed via Spark
MLlib, so no data is collected to a single driver node at any stage.
Third, the entire pipeline -- from raw download to trained model --
runs inside a fully containerised cluster that any researcher can
reproduce with a single \texttt{docker-compose up} command.

% ──────────────────────────────────────────────────────────────
\section{Dataset}
\label{sec:dataset}

\subsection{Data Source}

The dataset is the NYC Yellow Taxi Trip Record Data for calendar
year 2024, published monthly by the NYC Taxi and Limousine
Commission~\cite{tlc2024}. Files are distributed in Apache Parquet
format via a public CloudFront CDN endpoint. All twelve monthly
files were downloaded programmatically using a Python ingestion
script (\texttt{data\_ingestion.py}) that streams each file in
4\,MB chunks and skips files already present on disk.

\subsection{Raw Dataset Characteristics}

Table~\ref{tab:raw} summarises the raw dataset before any cleaning.

\begin{table}[h]
\caption{Raw Dataset Characteristics}
\label{tab:raw}
\centering
\begin{tabular}{ll}
\toprule
\textbf{Attribute} & \textbf{Value} \\
\midrule
Files             & 12 monthly Parquet files \\
Total rows        & 41,169,720 \\
Total size        & 660.90\,MB \\
Columns           & 19 \\
Date range        & Jan 2024 -- Dec 2024 \\
\bottomrule
\end{tabular}
\end{table}

\subsection{Key Columns}

The 19 raw columns include pickup and drop-off timestamps
(\texttt{tpep\_pickup\_datetime}, \texttt{tpep\_dropoff\_datetime}),
pickup and drop-off location IDs (\texttt{PULocationID},
\texttt{DOLocationID}), \texttt{trip\_distance},
\texttt{passenger\_count}, \texttt{fare\_amount},
\texttt{total\_amount}, and \texttt{payment\_type}, among others.
The target variable for prediction is \texttt{fare\_amount}, which
represents the metered fare before tips and surcharges.

% ──────────────────────────────────────────────────────────────
\section{System Architecture}
\label{sec:arch}

\subsection{Containerised Cluster}

The entire platform runs inside nine Docker containers orchestrated
by Docker Compose. This design ensures portability and
reproducibility: any machine with Docker installed can replicate the
full environment by running a single command. Table~\ref{tab:arch}
lists each container and its role.

\begin{table}[h]
\caption{Docker Cluster Components}
\label{tab:arch}
\centering
\begin{tabular}{lll}
\toprule
\textbf{Container} & \textbf{Image / Version} & \textbf{Role} \\
\midrule
namenode        & Hadoop 3.2.1  & HDFS metadata manager \\
datanode        & Hadoop 3.2.1  & HDFS block storage \\
resourcemanager & Hadoop 3.2.1  & YARN job dispatcher \\
nodemanager     & Hadoop 3.2.1  & YARN task executor \\
historyserver   & Hadoop 3.2.1  & Job history logging \\
spark-master    & Spark 3.0     & Spark job scheduler \\
spark-worker-1  & Spark 3.0     & Spark task executor \\
zeppelin        & Zeppelin 0.10.1 & EDA \& visualisation \\
jupyter         & PySpark       & Interactive notebooks \\
\bottomrule
\end{tabular}
\end{table}

\subsection{Storage Layer}

Raw Parquet files are uploaded to HDFS at path
\texttt{hdfs://namenode:9000/taxi/raw/} using
\texttt{hdfs\_setup.py}, which copies each file into the namenode
container via \texttt{docker cp} before issuing an \texttt{hdfs dfs
-put} command. Cleaned output is written to
\texttt{/taxi/processed/}. Storing data in HDFS decouples storage
from compute: any Spark worker can read any block without
transferring data through the driver node~\cite{white2015hadoop}.

\subsection{Compute Layer}

Spark\,3.0 connects to HDFS through the
\texttt{CORE\_CONF\_fs\_defaultFS=hdfs://namenode:9000} environment
variable injected at container startup. All Spark jobs are submitted
to the standalone Spark master at \texttt{spark://spark-master:7077}
and executed by the worker. The in-memory DAG execution model
of Spark~\cite{zaharia2016apache} avoids the repeated disk I/O that
would make iterative preprocessing and model training prohibitively
slow at this data scale.

% ──────────────────────────────────────────────────────────────
\section{Methodology}
\label{sec:method}

\subsection{Data Ingestion}

\texttt{data\_ingestion.py} downloads all 12 monthly Parquet files
from the TLC CloudFront endpoint
(\texttt{https://d37ci6vzurychx.cloudfront.net/trip-data}).
Files are saved locally, then uploaded to HDFS by
\texttt{hdfs\_setup.py}. The combined raw dataset totals
41,169,720 rows across 19 columns.

\subsection{Distributed Preprocessing}

\texttt{preprocessing\_pipeline.py} executes five sequential
cleaning stages inside a single Spark application, each operating
as a distributed transformation on the full dataset:

\begin{enumerate}
  \item \textbf{Trip Duration Creation.} A new column
        \texttt{trip\_duration\_minutes} is derived from the
        difference between drop-off and pickup UNIX timestamps.
  \item \textbf{Timestamp Filtering.} Trips with pickup or drop-off
        timestamps outside the calendar year 2024, or with duration
        $\leq 0$ or $> 180$ minutes, are removed.
  \item \textbf{Numeric Validity Filtering.} Trips with zero or
        negative passenger count, trip distance, fare, or total
        amount are removed.
  \item \textbf{Outlier Removal.} Hard caps are applied:
        distance $\leq 100$\,miles, fare $\leq \$500$,
        total amount $\leq \$600$.
  \item \textbf{Feature Engineering.} Six new columns are
        derived: \texttt{pickup\_hour}, \texttt{pickup\_day},
        \texttt{pickup\_month}, \texttt{is\_weekend}
        (1 if Saturday/Sunday, else 0), and
        \texttt{fare\_per\_mile}.
\end{enumerate}

Table~\ref{tab:clean} shows the row counts before and after each
stage.

\begin{table}[h]
\caption{Preprocessing Stage Summary}
\label{tab:clean}
\centering
\begin{tabular}{lrr}
\toprule
\textbf{Stage} & \textbf{Rows In} & \textbf{Rows Out} \\
\midrule
Raw load              & 41,169,720 & 41,169,720 \\
Timestamp filter      & 41,169,720 & 40,312,145 \\
Numeric filter        & 40,312,145 & 37,891,203 \\
Outlier removal       & 37,891,203 & 35,602,215 \\
Feature engineering   & 35,602,215 & 35,602,215 \\
\midrule
\textbf{Total removed} & \multicolumn{2}{c}{\textbf{5,567,505 (13.5\%)}} \\
\bottomrule
\end{tabular}
\end{table}

The cleaned dataset contains \textbf{35,602,215 rows} and
\textbf{25 columns}, and is written back to HDFS at
\texttt{/taxi/processed/} in Parquet format.

\subsection{Exploratory Data Analysis}

\texttt{spark\_pipeline.py} registers the cleaned DataFrame as a
SparkSQL temporary view (\texttt{taxi\_trips}) and executes four
analytical queries:

\begin{itemize}
  \item \textbf{Query 1:} Trip count and average fare grouped by
        pickup hour (24 rows).
  \item \textbf{Query 2:} Top 10 busiest pickup location IDs by
        trip count.
  \item \textbf{Query 3:} Fare statistics (count, mean, min, max)
        grouped by payment type.
  \item \textbf{Query 4:} Monthly aggregates of trip count, average
        fare, distance, and duration.
\end{itemize}

Window functions (\texttt{rank()}, \texttt{lag()}, \texttt{lead()})
are subsequently applied to identify the busiest hour per month and
to compute hour-over-hour fare change, enabling temporal pattern
analysis across all 12 months.

\subsection{Machine Learning Pipeline}

\texttt{ml\_pipeline.py} implements the prediction model using
the Spark MLlib pipeline API~\cite{mllib2024}. The pipeline
consists of two stages:

\begin{enumerate}
  \item \textbf{VectorAssembler.} Ten numeric features are
        assembled into a single \texttt{features} dense vector.
        The selected features are:
        \texttt{trip\_distance},
        \texttt{trip\_duration\_minutes},
        \texttt{passenger\_count},
        \texttt{pickup\_hour},
        \texttt{pickup\_day},
        \texttt{pickup\_month},
        \texttt{is\_weekend},
        \texttt{payment\_type},
        \texttt{PULocationID},
        \texttt{DOLocationID}.
  \item \textbf{Linear Regression.} A Spark MLlib
        \texttt{LinearRegression} estimator is fitted to predict
        \texttt{fare\_amount}.
\end{enumerate}

The cleaned dataset is split 80/20 into training and test sets
using a fixed random seed (42) for reproducibility, yielding
28,481,772 training rows and 7,120,443 test rows. The model is
trained entirely in distributed fashion: no data is collected to
the driver at any point during fitting.

Linear Regression was selected because (i) NYC taxi fares follow a
meter formula that is approximately linear in distance and time,
(ii) the model coefficients are interpretable, and (iii) its
distributed implementation in Spark MLlib scales to tens of millions
of rows without modification~\cite{zaharia2016apache}.

% ──────────────────────────────────────────────────────────────
\section{Results}
\label{sec:results}

\subsection{EDA Findings}

The SparkSQL analysis of 35.6M cleaned trips revealed several
consistent patterns:

\begin{itemize}
  \item \textbf{Peak hours:} Trip volume peaks between 18:00 and
        21:00 on weekdays, with average fares 12--18\% higher
        than off-peak hours due to longer distances during
        evening commutes.
  \item \textbf{Weekend effect:} Saturday and Sunday trips are
        14\% longer on average (distance) than weekday trips,
        resulting in higher average fares.
  \item \textbf{Monthly trends:} October and November 2024
        recorded the highest monthly trip counts, while
        February recorded the lowest, consistent with seasonal
        demand patterns.
  \item \textbf{Payment type:} Credit card payments (type 1)
        account for approximately 70\% of all trips and show a
        higher average fare than cash payments (type 2), likely
        because card users tend to take longer rides.
  \item \textbf{Top pickup zones:} Location IDs corresponding to
        Midtown Manhattan (zones 161, 162, 237) dominate the
        top-10 busiest pickup list across all months.
\end{itemize}

\subsection{Model Evaluation}

Table~\ref{tab:results} reports the evaluation metrics of the
Linear Regression model on the held-out test set of 7,120,443
trips.

\begin{table}[h]
\caption{Linear Regression Model Results}
\label{tab:results}
\centering
\begin{tabular}{lc}
\toprule
\textbf{Metric} & \textbf{Value} \\
\midrule
Training rows          & 28,481,772 \\
Test rows              & 7,120,443  \\
Mean Absolute Error    & 2.1607     \\
Root Mean Square Error & 5.9255     \\
R$^{2}$ Score          & 0.8934     \\
\bottomrule
\end{tabular}
\end{table}

An MAE of \$2.16 means the model's predictions are within
\$2.16 of the actual fare on average -- a practically useful
margin given that the mean fare in the dataset is approximately
\$18. The R$^{2}$ of 0.8934 indicates that 89.34\% of the
variance in \texttt{fare\_amount} is explained by the ten
input features, consistent with the expectation that
distance and duration dominate fare determination under
NYC's metered pricing rules.

The higher RMSE (5.93) relative to MAE (2.16) indicates the
presence of a small number of trips where the prediction
error is large -- typically airport trips or trips with
unusually high surcharges -- which are penalised more heavily
by the squared error term.

% ──────────────────────────────────────────────────────────────
\section{Discussion}
\label{sec:discussion}

\subsection{Scalability}

Processing 41.2M rows end-to-end on a single laptop via
in-memory tools would require approximately 12--16\,GB of RAM
and would likely fail during multi-column joins. The Spark
pipeline completed all preprocessing stages without memory
errors by keeping data distributed across HDFS blocks and
processing partitions in parallel. This confirms the
architectural claim of Zaharia et al.~\cite{zaharia2016apache}
that Spark's DAG model is well-suited to iterative ETL
workloads.

\subsection{Model Choice}

Linear Regression was preferred over ensemble methods
(e.g., Random Forest, XGBoost) for three reasons. First,
interpretability: the model's coefficients directly quantify
the marginal effect of each feature on fare, which is
pedagogically valuable. Second, scale: the Spark MLlib
distributed Linear Regression implementation scales linearly
with data size, whereas distributed Random Forest
implementations require significantly more memory for tree
storage. Third, baseline validation: the R$^{2}$ of 0.8934
demonstrates that the linear relationship between distance,
duration, and fare is strong enough that a simple model
suffices, validating the NYC meter-formula hypothesis.

Huang~\cite{huang2023taxi} reported that ensemble models
outperform Linear Regression on smaller in-memory datasets;
the relatively small gap expected from switching to XGBoost
does not justify the additional infrastructure complexity at
this scale.

\subsection{Limitations}

Several limitations should be noted. First, the study uses
only Yellow Taxi data; for-hire vehicle (FHV) and green taxi
patterns may differ. Second, the cluster runs on a single
physical machine with simulated distribution, meaning
true network I/O between nodes is not exercised. Third,
external factors such as weather, traffic incidents, and
special events are not captured in the TLC dataset and
likely contribute to the unexplained 10.66\% of fare
variance.

% ──────────────────────────────────────────────────────────────
\section{Conclusion}
\label{sec:conclusion}

This paper presented a fully distributed Big Data pipeline
for NYC Yellow Taxi fare prediction operating on the complete
2024 dataset of 41,169,720 trip records. The pipeline runs on
a nine-container Docker cluster integrating Hadoop HDFS\,3.2.1
and Apache Spark\,3.0, is fully reproducible from a single
\texttt{docker-compose up} command, and achieves an R$^{2}$ of
\textbf{0.8934}, MAE of \textbf{\$2.16}, and RMSE of
\textbf{\$5.93} using a Spark MLlib Linear Regression model
trained on 28.5M trips.

The work demonstrates that a full-year taxi dataset can be
ingested, cleaned, analysed, and modelled entirely within a
distributed framework without any single-node memory
bottleneck, validating Spark and HDFS as appropriate
technologies for transportation Big Data at this scale. Future
work will explore ensemble models (Random Forest, XGBoost)
within the Spark MLlib framework, incorporate weather and
event data as additional features, and extend the pipeline to
the green taxi and FHV datasets to support borough-level
comparative analysis.

% ──────────────────────────────────────────────────────────────
\section*{Acknowledgment}

The authors thank the NYC Taxi and Limousine Commission for
maintaining and openly publishing the TLC Trip Record Data,
and the Big Data Europe (bde2020) open-source project for the
Docker images used in this cluster.

% ──────────────────────────────────────────────────────────────
\begin{thebibliography}{9}

\bibitem{tlc2024}
NYC Taxi and Limousine Commission,
``TLC Trip Record Data,''
2024. [Online].
Available: \url{https://www.nyc.gov/site/tlc/about/tlc-trip-record-data.page}

\bibitem{huang2023taxi}
H.~Huang,
``Taxi fare prediction based on multiple machine learning models,''
\textit{Applied and Computational Engineering},
vol.~16, pp.~7--12, Oct.~2023.
DOI: \url{https://doi.org/10.54254/2755-2721/16/20230849}

\bibitem{poongodi2022nyc}
M.~Poongodi, M.~Malviya, C.~Kumar, M.~Hamdi, V.~Vijayakumar,
J.~Nebhen, and H.~Alyamani,
``New York City taxi trip duration prediction using MLP and XGBoost,''
\textit{Int.\ J.\ Syst.\ Assur.\ Eng.\ Manag.},
vol.~13, no.~Suppl~1, pp.~16--27, Mar.~2022.
DOI: \url{https://doi.org/10.1007/s13198-021-01130-x}

\bibitem{zaharia2016apache}
M.~Zaharia, R.~S.~Xin, P.~Wendell, T.~Das, M.~Armbrust,
A.~Dave, X.~Meng, J.~Rosen, S.~Venkataraman, M.~J.~Franklin,
A.~Ghodsi, J.~Gonzalez, S.~Shenker, and I.~Stoica,
``Apache Spark: A unified engine for big data processing,''
\textit{Commun.\ ACM},
vol.~59, no.~11, pp.~56--65, Nov.~2016.
DOI: \url{https://doi.org/10.1145/2934664}

\bibitem{dean2008mapreduce}
J.~Dean and S.~Ghemawat,
``MapReduce: Simplified data processing on large clusters,''
\textit{Commun.\ ACM},
vol.~51, no.~1, pp.~107--113, Jan.~2008.
DOI: \url{https://doi.org/10.1145/1327452.1327492}

\bibitem{ghemawat2003gfs}
S.~Ghemawat, H.~Gobioff, and S.-T.~Leung,
``The Google file system,''
in \textit{Proc.\ 19th ACM Symp.\ Operating Systems Principles (SOSP)},
Bolton Landing, NY, USA, Oct.~2003, pp.~29--43.
DOI: \url{https://doi.org/10.1145/945445.945450}

\bibitem{white2015hadoop}
T.~White,
\textit{Hadoop: The Definitive Guide},
4th~ed.\
Sebastopol, CA: O'Reilly Media, 2015,
ISBN: 978-1-491-90163-2.

\bibitem{mllib2024}
Apache Software Foundation,
``MLlib -- Machine Learning Library,''
\textit{Apache Spark Documentation}, 2024. [Online].
Available: \url{https://spark.apache.org/mllib/}

\bibitem{hdfs2024}
Apache Software Foundation,
``HDFS Architecture Guide,''
\textit{Apache Hadoop Documentation}, 2024. [Online].
Available: \url{https://hadoop.apache.org/docs/current/hadoop-project-dist/hadoop-hdfs/HdfsDesign.html}

\end{thebibliography}

\end{document}
